% =============================================================================
% Release 2 — Sprint 3 : Monitoring et supervision de la plateforme
% =============================================================================

\section[Sprint 3 : Monitoring et supervision de la plateforme]%
{Sprint 3 : Monitoring et supervision de la plateforme}

Ce sprint finalise la Release~2 par le \textbf{déploiement continu} (CD) de l'application
sur Kubernetes et la mise en place de la \textbf{supervision} (Prometheus, Grafana) via
Ansible sur vm-gestion.


% ---------------------------------------------------------------------------
\subsection{Objectif du Sprint 3}
% ---------------------------------------------------------------------------

Le tableau~\ref{tab:r2-s3-objectifs} détaille les objectifs du Sprint~3 par axe.

\begin{table}[H]
\centering
\caption{Objectifs du Sprint~3 par axe}
\label{tab:r2-s3-objectifs}
\footnotesize
\renewcommand{\arraystretch}{1.1}
\setlength{\tabcolsep}{6pt}
\begin{tabular}{|p{2.8cm}|p{9.2cm}|}
\hline
\textbf{Axe} & \textbf{Objectif} \\
\hline
Déploiement continu &
  Déployer automatiquement backend, frontend et MongoDB in-cluster via le stage
  \texttt{DeployK8s} (branche \texttt{main}, approbation manuelle) \\
\hline
Manifests Kubernetes &
  Appliquer les fichiers YAML du répertoire \texttt{k8s/} (Deployments, Services,
  StatefulSet MongoDB, Secrets) \\
\hline
Registre privé &
  Configurer le pull d'images depuis Nexus (\texttt{10.1.1.4:8082}) sur les workers \\
\hline
Supervision &
  Déployer Prometheus et Grafana sur vm-gestion (\texttt{ansible/playbooks/monitoring.yml}) \\
\hline
Observabilité &
  Collecter les métriques des nœuds K8s, du backend NestJS et de MongoDB \\
\hline
Validation &
  Smoke tests post-déploiement (API, frontend, cibles Prometheus UP) \\
\hline
\end{tabular}
\end{table}

\paragraph{Résultat attendu.}
\begin{itemize}
  \item namespace \texttt{ticketmanager} avec pods \texttt{Running} ;
  \item application accessible via NodePort (frontend \texttt{30080}, backend \texttt{30081}) ;
  \item Prometheus accessible sur le port \texttt{9090}, Grafana sur \texttt{3030} ;
  \item dashboards opérationnels pour le cluster et l'application.
\end{itemize}

\paragraph{Choix architectural — MongoDB in-cluster.}
Contrairement aux projets précédents (chatbot) où MongoDB était exclusivement sur une VM
dédiée, TicketManager héberge la base de données de production dans un
\textbf{StatefulSet Kubernetes} avec volume persistant. La vm-db reste disponible pour
migration ou sauvegarde legacy (\texttt{ansible/playbooks/mongodb.yml}).


% ---------------------------------------------------------------------------
\subsection{Backlog du Sprint 3}
% ---------------------------------------------------------------------------

Le tableau~\ref{tab:r2-s3-backlog-deploy} couvre l'epic déploiement continu. Le
tableau~\ref{tab:r2-s3-backlog-monitoring} couvre l'epic supervision.

\subsubsection{Epic — Déploiement continu (Kubernetes)}

\begin{table}[H]
\centering
\caption{Backlog Sprint~3 — Epic déploiement continu Kubernetes}
\label{tab:r2-s3-backlog-deploy}
\footnotesize
\renewcommand{\arraystretch}{0.85}
\setlength{\tabcolsep}{4pt}
\begin{tabular}{|c|p{7.8cm}|c|c|}
\hline
\textbf{ID} & \textbf{User story} & \textbf{Priorité} & \textbf{Est. (j)} \\
\hline
R2-S3-D01 &
  En tant que développeur, je déploie le backend NestJS en pod Kubernetes &
  H & 2 \\
\hline
R2-S3-D02 &
  En tant que développeur, je déploie le frontend Next.js en pod Kubernetes &
  H & 2 \\
\hline
R2-S3-D03 &
  En tant que développeur, je déploie MongoDB in-cluster via un StatefulSet &
  H & 2 \\
\hline
R2-S3-D04 &
  En tant que développeur, je configure les secrets Kubernetes (JWT, MongoDB, Cloudinary) &
  H & 1 \\
\hline
R2-S3-D05 &
  En tant que développeur, j'automatise le déploiement via le stage CD du pipeline &
  H & 2 \\
\hline
R2-S3-D06 &
  En tant que développeur, je valide le déploiement par des smoke tests automatisés &
  H & 1 \\
\hline
\end{tabular}
\end{table}

\subsubsection{Epic — Monitoring et supervision}

\begin{table}[H]
\centering
\caption{Backlog Sprint~3 — Epic monitoring et supervision}
\label{tab:r2-s3-backlog-monitoring}
\footnotesize
\renewcommand{\arraystretch}{0.85}
\setlength{\tabcolsep}{4pt}
\begin{tabular}{|c|p{7.8cm}|c|c|}
\hline
\textbf{ID} & \textbf{User story} & \textbf{Priorité} & \textbf{Est. (j)} \\
\hline
R2-S3-M01 &
  En tant que développeur, je déploie Prometheus sur vm-gestion via Ansible &
  H & 2 \\
\hline
R2-S3-M02 &
  En tant que développeur, je déploie Grafana et je connecte Prometheus comme source de données &
  H & 2 \\
\hline
R2-S3-M03 &
  En tant que développeur, j'intègre le déploiement monitoring au pipeline Azure DevOps &
  M & 1 \\
\hline
R2-S3-M04 &
  En tant que développeur, je vérifie que les cibles Prometheus sont en état \texttt{UP} &
  H & 1 \\
\hline
\end{tabular}
\end{table}


% ---------------------------------------------------------------------------
\subsection{Analyse — Conception du déploiement Kubernetes}
% ---------------------------------------------------------------------------

\subsubsection{Architecture de déploiement}

Les composants applicatifs sont encapsulés dans des pods du namespace
\texttt{ticketmanager} :
\begin{itemize}
  \item \textbf{ticket-backend} — image Nexus, port 3001, secrets JWT et Cloudinary ;
  \item \textbf{ticket-frontend} — image Nexus, port 3000, URL API figée au build Docker ;
  \item \textbf{mongodb} — StatefulSet avec volume persistant pour les données tickets.
\end{itemize}

Les requêtes externes atteignent les services via \textbf{NodePort} sur les workers
(\texttt{30080} frontend, \texttt{30081} backend). La figure~\ref{fig:r2-s3-k8s-arch}
détaille l'architecture de déploiement Kubernetes (ressources du namespace
\texttt{ticketmanager}), en complément de la vue d'ensemble physique (figure~2.5).

\begin{figure}[htbp]
    \centering
    \includegraphics[width=0.92\textwidth]{architecture-deploiement-k8s.png}
    \caption{Architecture de déploiement Kubernetes — namespace \texttt{ticketmanager}}
    \label{fig:r2-s3-k8s-arch}
\end{figure}

\subsubsection{Chaîne de déploiement continu}

Le stage \texttt{DeployK8s} du pipeline :
\begin{enumerate}
  \item attend une \textbf{approbation manuelle} (\texttt{ManualValidation}) ;
  \item configure le registre insecure sur les workers ;
  \item exécute \texttt{k8s/deploy-pipeline.sh} (secrets, manifests, rollout, seed, smoke test).
\end{enumerate}


% ---------------------------------------------------------------------------
\subsection{Réalisation — Déploiement Kubernetes}
% ---------------------------------------------------------------------------

\subsubsection{Manifests et secrets}

La figure~\ref{fig:r2-s3-backend-yaml} présente le manifest de déploiement du backend
(\texttt{backend-deployment.yaml}). La figure~\ref{fig:r2-s3-frontend-yaml} illustre celui
du frontend (\texttt{frontend-deployment.yaml}).

\begin{figure}[htbp]
    \centering
    \includegraphics[width=0.85\textwidth]{s3-backend-deployment.png}
    \caption{Manifest de déploiement du backend (\texttt{backend-deployment.yaml})}
    \label{fig:r2-s3-backend-yaml}
\end{figure}

\begin{figure}[htbp]
    \centering
    \includegraphics[width=0.85\textwidth]{s3-frontend-deployment.png}
    \caption{Manifest de déploiement du frontend (\texttt{frontend-deployment.yaml})}
    \label{fig:r2-s3-frontend-yaml}
\end{figure}

\subsubsection{Exécution du déploiement}

La figure~\ref{fig:r2-s3-k8s-pipeline} présente le stage \textbf{DeployK8s} exécuté dans
Azure DevOps, incluant l'approbation manuelle et l'application des manifests.

\begin{figure}[htbp]
    \centering
    \includegraphics[width=0.85\textwidth]{s3-k8s-pipeline-stage.png}
    \caption{Stage DeployK8s dans Azure DevOps}
    \label{fig:r2-s3-k8s-pipeline}
\end{figure}

\subsubsection{Vérification des pods}

La figure~\ref{fig:r2-s3-pods} présente la sortie de \texttt{kubectl get pods -n ticketmanager}
: l'ensemble des pods applicatifs est en état \texttt{Running}.

\begin{figure}[htbp]
    \centering
    \includegraphics[width=0.85\textwidth]{s3-kubectl-pods.png}
    \caption{État des pods (\texttt{kubectl get pods -n ticketmanager})}
    \label{fig:r2-s3-pods}
\end{figure}

\subsubsection{Application accessible}

La figure~\ref{fig:r2-s3-app} illustre l'interface TicketManager accessible depuis un
navigateur après déploiement réussi sur le cluster.

\begin{figure}[htbp]
    \centering
    \includegraphics[width=0.85\textwidth]{s3-app-browser.png}
    \caption{Interface TicketManager accessible après déploiement}
    \label{fig:r2-s3-app}
\end{figure}


% ---------------------------------------------------------------------------
\subsection{Réalisation — Monitoring et supervision}
% ---------------------------------------------------------------------------

\subsubsection{Déploiement Ansible (Prometheus + Grafana)}

Le playbook \texttt{ansible/playbooks/monitoring.yml} déploie Prometheus (port 9090) et
Grafana (port 3030) sur vm-gestion. La figure~\ref{fig:r2-s3-ansible} présente l'exécution
réussie de ce playbook.

\begin{figure}[htbp]
    \centering
    \includegraphics[width=0.85\textwidth]{s3-ansible-monitoring.png}
    \caption{Exécution du playbook \texttt{monitoring.yml}}
    \label{fig:r2-s3-ansible}
\end{figure}

\subsubsection{Configuration Prometheus}

Le fichier \texttt{ansible/templates/prometheus.yml.j2} définit les cibles surveillées.
La figure~\ref{fig:r2-s3-prom-targets} présente l'interface Prometheus (\texttt{/targets})
avec l'ensemble des cibles en état \texttt{UP}.

\begin{figure}[htbp]
    \centering
    \includegraphics[width=0.85\textwidth]{s3-prometheus-targets.png}
    \caption{Cibles actives dans l'interface Prometheus (\texttt{/targets})}
    \label{fig:r2-s3-prom-targets}
\end{figure}

\subsubsection{Dashboards Grafana}

La figure~\ref{fig:r2-s3-grafana-ds} présente la configuration de Prometheus comme source
de données dans Grafana. La figure~\ref{fig:r2-s3-grafana-dash} illustre un dashboard de
supervision regroupant les métriques du cluster et de l'application backend.

\begin{figure}[htbp]
    \centering
    \includegraphics[width=0.85\textwidth]{s3-grafana-datasource.png}
    \caption{Configuration de Prometheus comme source Grafana}
    \label{fig:r2-s3-grafana-ds}
\end{figure}

\begin{figure}[htbp]
    \centering
    \includegraphics[width=0.85\textwidth]{s3-grafana-dashboard.png}
    \caption{Dashboard de supervision — cluster et application}
    \label{fig:r2-s3-grafana-dash}
\end{figure}

\subsubsection{Intégration pipeline}

La figure~\ref{fig:r2-s3-mon-pipeline} présente le stage \textbf{DeployMonitoring} du
pipeline Azure DevOps, exécutant automatiquement le playbook Ansible après un déploiement
K8s réussi.

\begin{figure}[htbp]
    \centering
    \includegraphics[width=0.85\textwidth]{s3-monitoring-pipeline.png}
    \caption{Stage DeployMonitoring dans Azure DevOps}
    \label{fig:r2-s3-mon-pipeline}
\end{figure}


% ---------------------------------------------------------------------------
\subsection{Tests et validation}
% ---------------------------------------------------------------------------

\subsubsection{Smoke tests post-déploiement}

Le script \texttt{k8s/post-deploy.sh} vérifie la disponibilité de l'API, du frontend et
des outils de supervision. Le tableau~\ref{tab:r2-s3-validation} résume les critères de
validation retenus pour le Sprint~3.

\begin{table}[H]
\centering
\caption{Critères de validation du Sprint~3}
\label{tab:r2-s3-validation}
\footnotesize
\renewcommand{\arraystretch}{1.1}
\setlength{\tabcolsep}{6pt}
\begin{tabular}{|l|l|l|}
\hline
\textbf{Composant} & \textbf{URL / commande} & \textbf{Critère de succès} \\
\hline
Backend & \texttt{curl http://<worker>:30081/} & HTTP 200 \\
\hline
Frontend & \texttt{http://<worker>:30080} & Page de connexion affichée \\
\hline
Pods & \texttt{kubectl get pods -n ticketmanager} & Tous \texttt{Running} \\
\hline
Prometheus & \texttt{curl http://10.1.1.4:9090/-/ready} & \texttt{Ready} \\
\hline
Grafana & \texttt{http://10.1.1.4:3030/login} & HTTP 200 \\
\hline
\end{tabular}
\end{table}


% ---------------------------------------------------------------------------
\subsection{Bilan du Sprint 3}
% ---------------------------------------------------------------------------

Le Sprint~3 clôt la Release~2 en rendant TicketManager déployable et observable en
environnement cloud. L'application tourne dans Kubernetes ; la supervision Prometheus/Grafana
permet de détecter les anomalies de charge, de disponibilité et de performance avant
intervention manuelle.
